\section{Finite-domain area laws and polynomial PEPS approximation}
\label{sec:finite_domain_area_law}

An entropy bound for every cut and an approximation of the entire vector are
different assertions. The first concerns arbitrary finite induced lattice
domains. The second asks for one compatible tensor representation on the
original open square. We state the two models separately, following
\cite{OpenAI2026AreaLaw,OpenAI2026PolynomialPEPS}.

\begin{definition}[Induced lattice domain]
    \label{def:area_law_domain}
    \lean{
        TNLean.PEPS.AreaLaw.Site,
        TNLean.PEPS.AreaLaw.Configuration,
        TNLean.PEPS.AreaLaw.StateSpace,
        TNLean.PEPS.AreaLaw.domainGraph,
        TNLean.PEPS.AreaLaw.graphNeighborhood,
        TNLean.PEPS.AreaLaw.configurationSplit,
        TNLean.PEPS.AreaLaw.card_configuration
    }
    \leanok
    Let $\Lambda\subset\mathbb Z^2$ be finite. Join sites at coordinate distance
    one, and let $d_\Lambda$ be distance in this induced graph, with value
    $+\infty$ between distinct components. For $X\subseteq\Lambda$ and an integer
    $r\ge0$, the closed graph neighborhood
    \begin{align}
        N_r^\Lambda(X) & =\{y\in\Lambda:\ d_\Lambda(x,y)\le r
        \text{ for some }x\in X\}
        \label{eq:area_law_graph_neighborhood}
    \end{align}
    consists of the sites joined to some site of $X$ by an induced walk of
    length at most $r$. At every site put $\mathbb C^q$,
    $q\ge1$. In configuration coordinates the physical Hilbert space is
    \begin{align}
        \mathcal H_\Lambda     & =\mathbb C^{\{\sigma:\Lambda\to\{0,\ldots,q-1\}\}},
        \label{eq:area_law_state_space}                                              \\
        \dim\mathcal H_\Lambda & =q^{|\Lambda|}.
        \label{eq:area_law_dimension}
    \end{align}
    Restriction to $A$ and $\Lambda\setminus A$ identifies the configuration set
    with the product of the two regional configuration sets. Empty domains,
    holes, and disconnected components are permitted.
\end{definition}

\begin{definition}[Finite-range interactions and boundary]
    \label{def:area_law_hamiltonian}
    \lean{
        TNLean.PEPS.AreaLaw.IsAdmissibleSupport,
        TNLean.PEPS.AreaLaw.AdmissibleSupport,
        TNLean.PEPS.AreaLaw.LocalHamiltonian,
        TNLean.PEPS.AreaLaw.LocalHamiltonian.operator,
        TNLean.PEPS.AreaLaw.LocalHamiltonian.operator_isHermitian,
        TNLean.PEPS.AreaLaw.LocalHamiltonian.zero,
        TNLean.PEPS.AreaLaw.localLift,
        TNLean.PEPS.AreaLaw.localLift_mem_supportedOperators,
        TNLean.PEPS.AreaLaw.edgeBoundary,
        TNLean.PEPS.AreaLaw.edgeBoundary_empty,
        TNLean.PEPS.AreaLaw.edgeBoundary_univ
    }
    \leanok
    \uses{def:area_law_domain}
    Fix an integer $R\ge0$ and $J>0$. For each nonempty
    $X\subseteq\Lambda$ with $\operatorname{diam}_{d_\Lambda}X\le R$, assign
    one Hermitian operator $h_X$ supported on $X$, with $\|h_X\|\le J$.
    The assigned operator may be zero, and acts as the identity on the
    complementary tensor factors. Define
    \begin{align}
        H                  & =\sum_{\substack{\varnothing\ne X\subseteq\Lambda \\
                \operatorname{diam}_{d_\Lambda}X\le R}}h_X,
        \label{eq:area_law_hamiltonian}                                        \\
        \partial_\Lambda A & =
        \{\{x,y\}:x\in A,\ y\in\Lambda\setminus A,\ d_\Lambda(x,y)=1\}.
        \label{eq:area_law_boundary}
    \end{align}
    The norm is the Euclidean operator norm. Boundary edges are unordered and
    counted once; empty and full cuts have empty boundary.
\end{definition}

\begin{theorem}[Interaction range and rectangular coordinates]
    \label{thm:area_law_domain_identifications}
    \lean{
        TNLean.PEPS.AreaLaw.exists_walk_length_le_iff_edist_le,
        TNLean.PEPS.AreaLaw.isAdmissibleSupport_zero_iff,
        TNLean.PEPS.AreaLaw.isAdmissibleSupport_singleton,
        TNLean.PEPS.AreaLaw.isAdmissibleSupport_pair_of_adj,
        TNLean.PEPS.AreaLaw.rectangleEmbedding,
        TNLean.PEPS.AreaLaw.rectangleEmbedding_injective,
        TNLean.PEPS.AreaLaw.rectangularDomain,
        TNLean.PEPS.AreaLaw.rectangleSiteEquiv,
        TNLean.PEPS.AreaLaw.rectangleSiteEquiv_val,
        TNLean.PEPS.AreaLaw.rectangleGraphIso,
        TNLean.PEPS.AreaLaw.rectangleConfigurationEquiv,
        TNLean.PEPS.AreaLaw.isAdmissibleSupport_rectangle_pair_of_adj
    }
    \leanok
    \uses{def:area_law_domain,def:area_law_hamiltonian}
    A pair has induced distance at most $R$ if and only if an induced walk of
    length at most $R$ joins it. Every singleton is an admissible support at
    every range, and at range zero the admissible supports are exactly the
    singletons. Adjacent pairs are admissible at range one.
    The nearest-neighbor graph on
    $\{0,\ldots,w-1\}\times\{0,\ldots,h-1\}$ is the induced graph of that
    integer rectangle, with the corresponding identification of configurations
    and range-one pair supports.
\end{theorem}
\begin{proof}\leanok
    The extended graph distance is the least walk length, or infinity when no
    walk exists. A zero-length walk has equal endpoints. A nonempty support of
    diameter zero therefore has one element. The coordinate embedding preserves
    horizontal and vertical nearest neighbors and is bijective onto the integer
    rectangle. Restricting configurations along this bijection gives the last
    identification.
\end{proof}

\begin{definition}[Ground vector and regional entropy]
    \label{def:area_law_ground_entropy}
    \lean{
        TNLean.PEPS.AreaLaw.IsGappedGroundState,
        TNLean.PEPS.AreaLaw.reducedState,
        TNLean.PEPS.AreaLaw.reducedState_isHermitian,
        TNLean.PEPS.AreaLaw.regionalEntropy,
        TNLean.PEPS.AreaLaw.UniformAreaLaw
    }
    \leanok
    \uses{def:area_law_hamiltonian}
    For a unit vector $\Omega$, a real $E_0$, and $\Delta>0$, the eigenvector and
    full-system projector-gap conditions are
    \begin{align}
        H\Omega & =E_0\Omega,
        \label{eq:area_law_ground_vector}                    \\
        H-E_0I  & \ge\Delta(I-|\Omega\rangle\langle\Omega|).
        \label{eq:area_law_projector_gap}
    \end{align}
    For $A\subseteq\Lambda$, put
    \begin{align}
        \rho_A      & =\operatorname{Tr}_{\Lambda\setminus A}|\Omega\rangle\langle\Omega|,
        \label{eq:area_law_reduced_state}                                                  \\
        S_\Omega(A) & =-\operatorname{Tr}\rho_A\log\rho_A.
        \label{eq:area_law_regional_entropy}
    \end{align}
    All logarithms are natural, with $0\log0=0$. The reduced matrix is Hermitian,
    including when its spectrum contains zero. No condition is imposed on the
    gap of a Hamiltonian restricted to a smaller region.
    The uniform boundary assertion requires a single finite constant for each
    $q,R,J,\Delta$, chosen before every domain, Hamiltonian, ground vector, and cut.
\end{definition}

\begin{theorem}[Uniform finite-domain area law]
    \label{thm:openai_uniform_area_law}
    \uses{def:area_law_domain,def:area_law_hamiltonian,def:area_law_ground_entropy}
    For fixed integers $q\ge1$, $R\ge0$ and reals $J,\Delta>0$, there is a finite
    constant $C=C(q,R,J,\Delta)$ such that every finite induced domain with the
    above interaction bounds and a unique unit ground vector, up to phase,
    satisfying the full-system projector gap obeys
    $S_\Omega(A)\le C|\partial_\Lambda A|$ for every $A\subseteq\Lambda$.
    The same constant applies to all domains, interaction families, and cuts.
    This is Theorem~1.1 of \cite{OpenAI2026AreaLaw}.
\end{theorem}
\begin{proof}
    \uses{def:area_law_template,def:area_law_two_family_partition}
    For disjoint $B,C\subseteq\Lambda$ write
    $I_\Omega(B:C)=S_\Omega(B)+S_\Omega(C)-S_\Omega(B\cup C)$, and let $Z$ be the
    set of endpoints of the edges in $\partial_\Lambda A$. Empty and full cuts
    have zero entropy, and so does every cut with $\partial_\Lambda A=\varnothing$.
    Put $b=|\partial_\Lambda A|>0$.

    The global gap first gives an exponent below volume: for a box $Q$ of side
    at most $r$, sufficiently far from $Z$,
    \begin{align}
        S_\Omega(A\cap Q) & \le Cr^{1+e_0},\qquad e_0<1.
        \label{eq:area_law_initial_box}
    \end{align}
    Spectral filtering replaces $H$ by positive quasi-local constraints that
    annihilate $\Omega$ and keep a global gap. A replica argument then compares
    two bounds on a filtered norm while a finite randomized scan moves the
    partition through a collar. Repeating the comparison improves the exponent
    in (\ref{eq:area_law_initial_box}) and gives, for a template $T$ at scale
    $n$ and $L_n=\lfloor n^{1-\varepsilon}\rfloor$,
    \begin{align}
        I_\Omega\bigl(A\cap T:\Lambda\setminus(A\cap T_{L_n})\bigr)
         & \le Cn^{1-\varepsilon}.
        \label{eq:area_law_collar_information}
    \end{align}
    Here $\varepsilon$, $\alpha$ and $\beta$ are the exponents of
    Theorem~\ref{thm:area_law_exponent_gaps}. Products of the positive
    constraints amplify this bound. For every $M\ge1$
    there is a radius $r\le C_Mn^{(1-\varepsilon)/\alpha}$ with
    $L_n+2r=o(n^\beta)$ such that
    \begin{align}
        I_\Omega(A\cap T:J)
         & \le C_Mn^{-M}
        \qquad\text{for every }
        J\subseteq\Lambda\setminus N_{2r}^\Lambda(A\cap T_{L_n}).
        \label{eq:area_law_amplified_information}
    \end{align}

    The planar construction partitions $A$ into a residual set $D$ and an
    ordered family of pieces $X_i\subseteq A\cap T_i$ with labels
    $f_i\in\{0,1\}$, templates $T_i$ at scales $n_i$, and
    \begin{align}
        |D| & \le Cb,
        \qquad
        \sum_in_i^{-100}\le Cb,
        \label{eq:area_law_partition_counts} \\
        \operatorname{dist}_\infty(T_i,Z)
            & >2n_i^\beta,
        \qquad
        \operatorname{dist}_\infty\Bigl(T_i,
        \bigcup_{j<i,\,f_j=f_i}X_j\Bigr)>2n_i^\beta.
        \label{eq:area_law_partition_clearance}
    \end{align}
    Put $S_i=A\cap(T_i)_{L_{n_i}}$ and
    $J_i=(\Lambda\setminus A)\cup\bigcup_{j<i,\,f_j=f_i}X_j$, and let $r_i$ be
    the radius in (\ref{eq:area_law_amplified_information}) at scale $n_i$, with
    the scale threshold large enough that $L_{n_i}+2r_i<n_i^\beta$. A walk of
    length at most $2r_i$ starting in $S_i$ stays within ambient distance
    $n_i^\beta$ of $T_i$. By (\ref{eq:area_law_partition_clearance}) it meets no
    earlier piece of the same family, and it cannot leave $A$, since its first
    step out of $A$ would end at a point of $Z$. Hence
    $J_i\cap N_{2r_i}^\Lambda(S_i)=\varnothing$. With $M=100$,
    (\ref{eq:area_law_amplified_information}) and monotonicity of mutual
    information give
    \begin{align}
        I_\Omega(X_i:J_i)\le I_\Omega(A\cap T_i:J_i)
         & \le Cn_i^{-100}=:\eta_i.
        \label{eq:area_law_piece_information}
    \end{align}
    For one family $f$, with $A_f=\bigcup_{f_i=f}X_i$ and
    $E_f=\sum_{f_i=f}\eta_i$, the chain rule, purity and subadditivity
    give
    \begin{align}
        S_\Omega(A)+\sum_{f_i=f}S_\Omega(X_i)-E_f
         & \le S_\Omega\bigl((\Lambda\setminus A)\cup A_f\bigr)
        =S_\Omega(D\cup A_{1-f})
        \nonumber                                               \\
         & \le S_\Omega(D)+\sum_{f_i=1-f}S_\Omega(X_i).
        \label{eq:area_law_family_chain}
    \end{align}
    Adding (\ref{eq:area_law_family_chain}) for $f=0$ and $f=1$ cancels every
    piece entropy. With (\ref{eq:area_law_partition_counts}),
    \begin{align}
        S_\Omega(A)\le S_\Omega(D)+\frac12\sum_i\eta_i
         & \le|D|\log q+\frac C2\sum_in_i^{-100}\le C'b.
        \label{eq:area_law_final_bound}
    \end{align}
    Every threshold and constant above depends only on $q,R,J,\Delta$, and is
    chosen before the domain and the cut.
\end{proof}

\begin{definition}[Open-square Hamiltonian and PEPS approximation]
    \label{def:polynomial_peps_model}
    \lean{
        TNLean.PEPS.Approximation.Configuration,
        TNLean.PEPS.Approximation.StateSpace,
        TNLean.PEPS.Approximation.SquareHamiltonian,
        TNLean.PEPS.Approximation.SquareHamiltonian.operator,
        TNLean.PEPS.Approximation.SquareHamiltonian.IsGappedGroundState,
        TNLean.PEPS.Approximation.pepsVector,
        TNLean.PEPS.Approximation.HasPEPSApproximation,
        TNLean.PEPS.Approximation.PolynomialPEPSApproximation
    }
    \leanok
    On the original open $L\times L$ graph, let
    \begin{align}
        H       & =\sum_v h_v+\sum_e h_e,
        \label{eq:polynomial_peps_hamiltonian} \\
        \|h_v\| & \le J,
        \label{eq:polynomial_peps_site_bound}  \\
        \|h_e\| & \le J,
        \label{eq:polynomial_peps_edge_bound}
    \end{align}
    where terms are supported on their designated sites or edges, and only the
    total $H$ is assumed Hermitian. For a unit vector $\Omega$, a real $E_0$,
    and $\Delta>0$, the gapped ground-vector condition for this $H$ is
    \begin{align}
        H\Omega=E_0\Omega,
        \qquad
        H-E_0I & \ge\Delta(I-|\Omega\rangle\langle\Omega|).
        \label{eq:polynomial_peps_projector_gap}
    \end{align}
    A PEPS approximation with parameters $C,c$
    is a tensor family on those same edges with positive dimensions $D_e$
    satisfying $D_e\le CL^c$ for every edge $e$, and a nonzero contracted vector
    $\Phi$ such that, for some $\theta\in\mathbb R$,
    \begin{align}
        \left\|\frac{\Phi}{\|\Phi\|}-e^{i\theta}\Omega\right\| & \le L^{-1}.
        \label{eq:polynomial_peps_vector_error}
    \end{align}
    The norm in (\ref{eq:polynomial_peps_vector_error}) is the Euclidean norm
    of the entire vector.
    The local tensor maps need not be isometries and the bond dimensions need not
    be equal. The uniform approximation assertion requires $C,c>0$ to be chosen
    for fixed $q\ge2$, $J,\Delta>0$ before all $L\ge2$, Hamiltonians, and ground vectors.
\end{definition}

\begin{theorem}[Polynomial PEPS approximation]
    \label{thm:openai_polynomial_peps}
    \uses{def:polynomial_peps_model}
    Fix an integer $q\ge2$ and reals $J,\Delta>0$. There are $C,c>0$, depending
    only on these parameters, such that for every $L\ge2$, every Hamiltonian in
    the preceding model with a unique unit ground vector, up to phase, satisfying
    (\ref{eq:polynomial_peps_projector_gap}) has a nonzero PEPS on the original open
    grid with $D_e\le CL^c$ for every edge $e$ and
    \begin{align}
        \min_{\theta\in\mathbb R}
        \left\|\frac{\Phi}{\|\Phi\|}-e^{i\theta}\Omega\right\| & \le L^{-1}.
        \label{eq:polynomial_peps_phase_error}
    \end{align}
    This is Theorem~1.1 of \cite{OpenAI2026PolynomialPEPS}.
\end{theorem}
\begin{proof}
    \uses{thm:openai_uniform_area_law}
    Write $\Lambda_L$ for the sites of the grid and $\mathcal H_X$ for the
    tensor product of the site spaces in $X\subseteq\Lambda_L$. The physical
    registers are distributed among parties indexed by dyadic squares, starting
    from one party holding $\Omega$ and ending with one party per site. The
    area-law and collar estimates make the mutual information across tapered
    buffers smaller than any fixed inverse power of $L$, which
    lets an owner split its purification approximately. Near a square of
    logarithmic size, the ground vector is projected with high accuracy onto a
    sum of mutually orthogonal subspaces
    \begin{align}
        \operatorname{span}\{v_{jk}\}_{k=1}^{d_j}
        \otimes\mathcal H_{\Lambda_L\setminus X_j},
        \qquad
        \sum_jd_j & \le C_0L^{c_0},
        \label{eq:polynomial_peps_rank_bound}
    \end{align}
    with nested square samples $X_j$ and orthonormal vectors
    $v_{jk}\in\mathcal H_{X_j}$. Each party takes part in boundedly many
    operations with other parties during its whole lifetime, and each such
    operation is a sum of polynomially many products of private contractions
    and bipartite vectors or covectors.

    Telescoping the accumulated errors of these contractions, the final vector
    $\Psi$ and the final reference vector $\Omega\otimes\Xi$ satisfy
    \begin{align}
        \|\Psi-\Omega\otimes\Xi\| & \le CL^{-17},
        \qquad\|\Psi\|\le1,
        \qquad 1-CL^{-57}\le\|\Xi\|\le1,
        \label{eq:polynomial_peps_protocol_error} \\
        \bigl\|\operatorname{Tr}_{\rm private}|\Psi\rangle\langle\Psi|
        -|\Omega\rangle\langle\Omega|\bigr\|_1
                                  & \le CL^{-10}.
        \label{eq:polynomial_peps_protocol_density}
    \end{align}
    Density compression replaces the bipartite resources by polynomial sums of
    product operators and gives an operator network $\sigma$ on the party graph
    with polynomial link dimensions and
    \begin{align}
        \bigl\|\sigma-|\Omega\rangle\langle\Omega|\bigr\|_1
         & \le\eta:=CL^{-5}.
        \label{eq:polynomial_peps_compressed_density}
    \end{align}
    If $\eta<1$, write $E=\sigma-|\Omega\rangle\langle\Omega|$ and
    $a_z=\langle\Omega|z\rangle$ for product basis vectors $|z\rangle$. Since
    $\sum_z\|E|z\rangle\|^2\le\eta^2$ and $\sum_z|a_z|^2=1$, some $z$ has
    $a_z\ne0$ and $\|E|z\rangle\|\le\eta|a_z|$. The column $v=\sigma|z\rangle$
    is then nonzero, has the same virtual dimensions, and satisfies
    \begin{align}
        \min_{\theta\in\mathbb R}
        \left\|\frac{v}{\|v\|}-e^{i\theta}\Omega\right\|
         & \le2\eta.
        \label{eq:polynomial_peps_column_error}
    \end{align}
    Routing each link horizontally and then vertically through the anchors of
    the dyadic blocks has edge congestion at most a constant $\chi$, also after
    folding the padded grid into the $L\times L$ grid. If every link has
    dimension at most $C_1L^k$, grouping the links crossing each grid edge
    gives
    \begin{align}
        D_e & \le(C_1L^k)^\chi=C_1^\chi L^{k\chi}.
        \label{eq:polynomial_peps_routed_bond}
    \end{align}
    For $L\ge L_0$ the error $2CL^{-5}$ is at most $L^{-1}$, so
    (\ref{eq:polynomial_peps_column_error}) gives
    (\ref{eq:polynomial_peps_phase_error}). Each of the finitely many sizes
    $2\le L<L_0$ has an exact representation with bond dimensions bounded
    independently of the Hamiltonian, which is absorbed by enlarging the
    prefactor in (\ref{eq:polynomial_peps_routed_bond}).
\end{proof}

\begin{definition}[Polygonal lattice templates]
    \label{def:area_law_template}
    \lean{
        TNLean.PEPS.AreaLaw.Geometry.IsAllowedSlope,
        TNLean.PEPS.AreaLaw.Geometry.integerPoint,
        TNLean.PEPS.AreaLaw.Geometry.TemplatePolygon,
        TNLean.PEPS.AreaLaw.Geometry.TemplatePolygon.region,
        TNLean.PEPS.AreaLaw.Geometry.Template,
        TNLean.PEPS.AreaLaw.Geometry.Template.IsSeparated,
        TNLean.PEPS.AreaLaw.Geometry.Template.core,
        TNLean.PEPS.AreaLaw.Geometry.Template.collar,
        TNLean.PEPS.AreaLaw.ambientDilation,
        TNLean.PEPS.AreaLaw.ambientDilation_zero
    }
    \leanok
    \uses{def:area_law_domain,def:area_law_hamiltonian}
    Fix $D_0>0$ and a sufficiently large positive numerical constant
    $C_{\rm tpl}$ before choosing templates. A template with positive integer
    parameters $n,s_0$ is a nonempty finite
    $T\subset\mathbb Z^2$ with
    \begin{align}
        T                            & =\bigcup_{a=0}^{N_0-1}(P_a\cap\mathbb Z^2),
        \label{eq:area_law_template_cover}                                         \\
        N_0                          & \ge1,
        \label{eq:area_law_template_pieces}                                        \\
        C_{\rm tpl}N_0(s_0+1)        & \le n,
        \label{eq:area_law_template_scale}                                         \\
        \operatorname{diam}_\infty T & \le n.
        \label{eq:area_law_template_diameter}
    \end{align}
    Each $P_a$ is a closed convex rectangle or triangle of sup-norm diameter at
    most $s_0$, with side slopes in $\{0,\infty,1,-1\}$. Triangles have
    noncollinear vertices and rectangles have nonzero sides; this ordinary
    reading of the source's polygons is recorded in
    \cite{gap:arealaw2d_nondegenerate_templates}. The union may be
    disconnected. Geometric estimates require a uniform lower bound on
    $C_{\rm tpl}$, chosen independently of the template.
    Separation from a fixed cut set $Z$ means
    $\operatorname{dist}_\infty(T,Z)>4D_0s_0$. For the fixed cut
    $A\subseteq\Lambda$ of (\ref{eq:area_law_boundary}), put
    \begin{align}
        T_j & =T+([-j,j]^2\cap\mathbb Z^2),
        \label{eq:area_law_template_dilation} \\
        X   & =A\cap T,
        \label{eq:area_law_template_core}     \\
        Q_j & =A\cap(T_j\setminus T).
        \label{eq:area_law_template_collar}
    \end{align}
    In particular, $T_0=T$. These are ambient dilations, not the induced-graph
    neighborhoods (\ref{eq:area_law_graph_neighborhood}).
    This is Definition~9.3 of \cite{OpenAI2026AreaLaw}.
\end{definition}

\begin{definition}[Ordered two-family partition]
    \label{def:area_law_two_family_partition}
    \lean{
        TNLean.PEPS.AreaLaw.Geometry.OrderedTwoFamilyPartition,
        TNLean.PEPS.AreaLaw.Geometry.OrderedTwoFamilyPartition.earlierSameFamily
    }
    \leanok
    A finite region $A$ is partitioned into a residual set $D$ and pairwise
    disjoint pieces $X_i$. Each piece has a family label in $\{0,1\}$, and one
    finite order induces the order within each family. For a piece $X_i$ in
    family $f$, the earlier same-family region is
    $\bigcup_{j<i,\,f(j)=f}X_j$. No size, separation, or information bound is
    part of this partition data; those are additional mathematical assertions.
\end{definition}

\begin{theorem}[Exact exponent gaps]
    \label{thm:area_law_exponent_gaps}
    \lean{
        TNLean.PEPS.AreaLaw.Exponents.boxError,
        TNLean.PEPS.AreaLaw.Exponents.scannerMu,
        TNLean.PEPS.AreaLaw.Exponents.scannerNu,
        TNLean.PEPS.AreaLaw.Exponents.scannerEll,
        TNLean.PEPS.AreaLaw.Exponents.scannerOmega,
        TNLean.PEPS.AreaLaw.Exponents.scannerKappa,
        TNLean.PEPS.AreaLaw.Exponents.amplificationEpsilon,
        TNLean.PEPS.AreaLaw.Exponents.beta,
        TNLean.PEPS.AreaLaw.Exponents.amplificationPower,
        TNLean.PEPS.AreaLaw.Exponents.alpha,
        TNLean.PEPS.AreaLaw.Exponents.geometryDelta,
        TNLean.PEPS.AreaLaw.Exponents.zeta,
        TNLean.PEPS.AreaLaw.Exponents.rho,
        TNLean.PEPS.AreaLaw.Exponents.geometryEll,
        TNLean.PEPS.AreaLaw.Exponents.alpha_gt,
        TNLean.PEPS.AreaLaw.Exponents.geometry_gaps
    }
    \leanok
    The final box error exponent is $e_*=2\cdot10^{-6}$, and the scanner choices are
    $(\mu,\nu,\ell_{\rm sc},\omega,\kappa)=(1/4,1/1000,1/100000,1/20000,1/1000000)$.
    The amplification uses $\varepsilon=10^{-5}$ and $\alpha=p/(p+1)$.
    Set $\beta=1-2\cdot10^{-6}$, $\delta_0=2\cdot10^{-7}$,
    $\zeta=1-\delta_0/2$, $\rho_0=(1+\beta)/2$, and $\ell_{\rm geo}=10^{-5}$.
    Then
    \begin{align}
        (1+2\delta_0)\beta              & <\zeta<1,
        \label{eq:area_law_zeta_gap}                       \\
        (1+2\delta_0)(1-\ell_{\rm geo}) & <1,
        \label{eq:area_law_ell_gap}                        \\
        1-\ell_{\rm geo}                & <\beta<\rho_0<1.
        \label{eq:area_law_rho_gap}
    \end{align}
    The choice $p=10^6$ also satisfies $p/(p+1)>1-10^{-6}$. The scanner
    exponent $\ell_{\rm sc}$ and the geometric exponent $\ell_{\rm geo}$ are
    separate choices, which happen to have the same value.
\end{theorem}
\begin{proof}\leanok
    Substitute the stated rational values and compare integers after multiplying
    by positive denominators. Each gap is strictly positive.
\end{proof}
